\documentclass{article}

% Language and encoding
\usepackage[portuguese]{babel}
\usepackage[utf8]{inputenc}
\usepackage[T1]{fontenc}

% Page layout
\usepackage[a4paper,left=25mm,right=25mm,top=25mm,bottom=25mm,headheight=6mm,footskip=12mm]{geometry}
\setlength{\parindent}{1em}
\setlength{\parskip}{1ex}

% Core packages
\usepackage{amsmath}
\usepackage{booktabs}
\usepackage{indentfirst}
\usepackage{changepage}
\usepackage{ragged2e}
\usepackage{setspace}
\usepackage{caption}
\usepackage{float}
\usepackage{subcaption}
\usepackage{longtable}
\usepackage{multirow}
\usepackage{pgfgantt}
\usepackage{titlesec}
\usepackage{multicol}
\usepackage{graphicx}
\usepackage{rotating}
\usepackage{adjustbox}
\usepackage{lmodern}
\usepackage{tocloft}
\usepackage{placeins}
\usepackage{xcolor}
\usepackage{listings}
\usepackage{lastpage}
\usepackage{fancyhdr}
\usepackage{hyperref}

%% headers & footers
\fancyhf{}
\rhead{\small{\emph{\projtitle, \projauthor}}}
\rfoot{\small{\thepage\ / \pageref{LastPage}}}
\pagestyle{fancy}
\renewcommand{\headrulewidth}{0.4pt}
\renewcommand{\footrulewidth}{0.4pt}

%% colors
\definecolor{engineering}{rgb}{0.549,0.176,0.098}
\definecolor{cloudwhite}{cmyk}{0,0,0,0.025}

%% source-code listings
\lstset{
 language=Java,
 basicstyle=\footnotesize\ttfamily,
 keywordstyle=\bfseries,
 numbers=left,
 numberstyle=\scriptsize\texttt,
 stepnumber=1,
 numbersep=8pt,
 frame=tb,
 float=htb,
 aboveskip=8mm,
 belowskip=4mm,
 backgroundcolor=\color{cloudwhite},
 showspaces=false,
 showstringspaces=false,
 showtabs=false,
 tabsize=2,
 captionpos=b,
 belowcaptionskip=12pt,
 breaklines=true,
 breakatwhitespace=false,
 commentstyle=\color{gray},
 mathescape=true,
 lineskip=3pt,
 extendedchars=true
}

%% hyperreferences
\hypersetup{
 plainpages=false,
 pdfpagelayout=SinglePage,
 bookmarksopen=false,
 bookmarksnumbered=true,
 breaklinks=true,
 linktocpage,
 colorlinks=true,
 linkcolor=engineering,
 urlcolor=engineering,
 filecolor=engineering,
 citecolor=engineering,
 allcolors=engineering
}

%% path to figures
\graphicspath{{figures/}}

%% macros
\newcommand{\school}{Universidade do Minho}
\newcommand{\degree}{Licenciatura em Engenharia Informática}
\newcommand{\projtitle}{Análise e Teste de Software}
\newcommand{\subtitle}{2025/2026}
\newcommand{\projauthor}{Trabalho Prático}
\newcommand{\projname}{Projetos 1 e 2 - SpotifyUM/SpotifUM}

\newcommand{\atsfigure}[3]{
\begin{figure}[H]
\centering
\IfFileExists{#1}{
    \includegraphics[width=#2\textwidth]{#1}
}{
    \fbox{\parbox[c][50mm][c]{0.9\textwidth}{\centering Placeholder para print: #1}}
}
\caption{#3}
\end{figure}
}

\renewcommand{\thesection}{\arabic{section}}
\setcounter{secnumdepth}{3}

\begin{document}

% Capa
\begin{titlepage}
\centering

\vspace{-15mm}
{\large \textbf{\textsc{\school}}}\\

\vspace{0.5cm}
{\large \degree}\\[8mm]

\vspace{2cm}
\includegraphics[width=52mm]{UM.png}

\vspace{1cm}
{\Large \projtitle}\\[2mm]
{\large \subtitle}\\[28mm]

{\Large \textbf{\projauthor}}\\
\vspace{0.5cm}
{\Large \projname}

\vspace{5cm}

{\large Francisco Barros A106894}\\[2mm]
{\large Tiago Martins A106927}\\[2mm]

\vspace{2cm}
{\Large \textbf{Maio 2026}}

\end{titlepage}

\renewcommand{\contentsname}{Índice}
\tableofcontents
\newpage
\listoffigures
\newpage
\listoftables
\newpage

\section{Introdução}
O presente relatório descreve a aplicação de técnicas de Análise e Teste de Software aos Projetos 1 e 2 (SpotifyUM/SpotifUM), no contexto da unidade curricular de ATS. O objetivo foi elevar a maturidade de teste das soluções através de uma abordagem integrada, combinando:
\begin{itemize}
    \item reforço manual de testes unitários com JUnit;
    \item avaliação estrutural por cobertura (JaCoCo);
    \item avaliação de qualidade dos testes por mutação (PIT);
    \item geração automática de propriedades com jqwik (Property-Based Testing);
    \item tentativa de geração automática de testes com EvoSuite;
    \item automação integral do pipeline para execução reprodutível.
\end{itemize}

Este documento apresenta não apenas os resultados obtidos, mas também as decisões técnicas, limitações da toolchain e o racional de melhoria iterativa usado para aumentar a robustez da suite de testes.

\section{Objetivos e Requisitos}
\subsection{Objetivos ATS}
\begin{itemize}
    \item Complementar e reforçar a suite JUnit.
    \item Executar geração automática com EvoSuite.
    \item Analisar cobertura e qualidade de testes.
    \item Avaliar robustez com mutação (PIT).
    \item Aplicar Property-Based Testing em larga escala.
\end{itemize}

\subsection{Extras}
\begin{itemize}
    \item Automação do pipeline com um comando.
    \item Refatoração e validação por regressão.
    \item Reforço de propriedades com PBT.
\end{itemize}

\subsection{Estado de cumprimento}
\begin{table}[H]
\centering
\begin{tabular}{p{0.52\textwidth}cc}
\toprule
Requisito do enunciado ATS & Projeto 1 & Projeto 2 \\
\midrule
Complemento da suite JUnit & Concluído & Concluído \\
EvoSuite para geração automática & N/A & Concluído com limitação de runtime \\
Análise de cobertura e qualidade (JaCoCo) & N/A & Concluído \\
Mutação com PIT & N/A & Concluído \\
PBT em larga escala (jqwik) & N/A & Concluído \\
Relatório com análise e gráficos & Concluído & Concluído (com placeholders de prints) \\
\bottomrule
\end{tabular}
\caption{Matriz de cumprimento dos requisitos do enunciado}
\end{table}

\section{Projeto 1 — SpotifyUM (Maven)}
\subsection{Objetivo}
O objetivo no Projeto 1 foi complementar a suite JUnit existente com testes unitários determinísticos em módulos que estavam sem cobertura (utilitários e controlador), privilegiando cenários de maior risco (persistência e fluxos base de autenticação).

\subsection{Testes adicionados}
Foram criadas as seguintes classes de teste (Projeto 1):
\begin{itemize}
    \item \texttt{Projeto1/SpotifyUM/src/test/java/org/Utils/SerializationTest.java}
    \item \texttt{Projeto1/SpotifyUM/src/test/java/org/Utils/MusicPlayerTest.java}
    \item \texttt{Projeto1/SpotifyUM/src/test/java/org/Utils/BrowserOpenerTest.java}
    \item \texttt{Projeto1/SpotifyUM/src/test/java/org/Controller/ControllerTest.java}
\end{itemize}

Principais cenários cobertos:
\begin{itemize}
    \item \textbf{Serialization}: round-trip exportar/importar com \texttt{@TempDir}, e validação de exceções (objeto nulo, caminho inválido, ficheiro inexistente).
    \item \textbf{MusicPlayer}: retorno \texttt{null} quando não existe WAV correspondente no diretório de recursos.
    \item \textbf{BrowserOpener}: validação de erro para URL inválido (\texttt{URISyntaxException}).
    \item \textbf{Controller}: mensagens e efeitos para registo duplicado, login com sucesso/insucesso, e import/export do modelo para ficheiro temporário.
\end{itemize}

\subsection{Execução e evidência}
A suite foi validada com Maven no Projeto 1:
\begin{lstlisting}[caption={Execução da suite de testes do Projeto 1 (Maven)},label={lst:proj1-mvn-test}]
cd Projeto1/SpotifyUM
mvn test
\end{lstlisting}

Resultado observado: execução concluída com sucesso (\texttt{exit=0}).

\subsection{EvoSuite no Projeto 1: integração, diagnóstico e estado atual}
Nesta iteração foi integrada a execução de EvoSuite no Projeto 1 (Maven), com foco em manter o escopo restrito a \texttt{Projeto1/SpotifyUM} e em registar as decisões técnicas tomadas.

\subsubsection{Alterações aplicadas}
Foram realizadas as seguintes alterações de configuração:
\begin{itemize}
    \item atualização de \texttt{pom.xml} com repositórios \texttt{https://www.evosuite.org/m2} em \texttt{repositories} e \texttt{pluginRepositories};
    \item definição de \texttt{evosuite.version=1.0.3} (versão efetivamente disponível no repositório EvoSuite);
    \item substituição da integração por plugin legado (que falhava na inicialização de classe) por execução \textit{standalone} do EvoSuite via \texttt{exec-maven-plugin};
    \item parametrização de \texttt{cuts} e \texttt{search\_budget} para execução incremental por classe;
    \item inclusão automática de \texttt{evosuite-tests} como fonte de testes com \texttt{build-helper-maven-plugin};
    \item adição de dependências \texttt{junit:junit} e runtime EvoSuite para compilar/executar os testes gerados;
    \item ajuste de compilação para Java 8 no Projeto 1, por compatibilidade do ecossistema EvoSuite 1.0.3;
    \item ajustes de compatibilidade Java 8 em código fonte/testes (substituição de APIs Java 9+ como \texttt{List.of}, \texttt{Optional.isEmpty}, \texttt{String.repeat}).
\end{itemize}

\subsubsection{Comandos usados na validação}
\begin{lstlisting}[caption={Validação de build normal e execução EvoSuite no Projeto 1},label={lst:proj1-evosuite-commands}]
# build normal do Projeto 1
mvn -f Projeto1/SpotifyUM/pom.xml clean test

# execução EvoSuite (JDK 8)
JAVA_HOME=/usr/lib/jvm/java-8-openjdk-amd64 \
PATH="/usr/lib/jvm/java-8-openjdk-amd64/bin:$PATH" \
mvn -f Projeto1/SpotifyUM/pom.xml -P evosuite \
  -Devosuite.cuts=org.Model.Music.Music \
  -Devosuite.searchBudgetInSeconds=30 \
  clean test
\end{lstlisting}

\subsubsection{Diagnóstico da falha de geração}
Na primeira integração (plugin legado), o EvoSuite reportou:
\begin{itemize}
    \item \texttt{WARN: failed to generate tests for <CUT>};
    \item logs em \texttt{.evosuite/.../std\_err\_CLIENT.log} com \texttt{Error while initializing target class: java.lang.IllegalArgumentException}.
\end{itemize}

Após migração para execução standalone (classpath explícito em \texttt{target/classes}), a geração passou a produzir classes \texttt{\_ESTest} em \texttt{evosuite-tests}.

\subsubsection{Estado atual}
\begin{itemize}
    \item \textbf{Concluído}: integração técnica do EvoSuite no Projeto 1 via profile Maven dedicado;
    \item \textbf{Concluído}: execução reprodutível do profile \texttt{evosuite};
    \item \textbf{Concluído}: geração efetiva de testes (ex.: \texttt{Music\_ESTest} e \texttt{Music\_ESTest\_scaffolding});
    \item \textbf{Concluído}: execução da suite após geração com \texttt{BUILD SUCCESS}.
\end{itemize}

\subsubsection{Resultados finais obtidos}
Na execução final validada para o Projeto 1, o EvoSuite gerou e executou com sucesso a suite da classe \texttt{org.Model.Music.Music}.

\begin{table}[H]
\centering
\begin{tabular}{lr}
\toprule
Métrica & Valor \\
\midrule
Classes testáveis alvo & 1 \\
Coverage global & 0.99 \\
Test suites geradas & 1 \\
Teste gerados pelo EvoSuite & 48 \\
Statements totais & 112 \\
Tempo de esforço & 37 s \\
Cobertura de linha & 100\% \\
Cobertura de branch & 100\% \\
Cobertura de exceções & 100\% \\
Weak mutation score & 100\% \\
Cobertura de output & 93\% \\
Cobertura de métodos & 100\% \\
Cobertura CBranch & 100\% \\
\bottomrule
\end{tabular}
\caption{Resultados finais do EvoSuite no Projeto 1}
\end{table}

\noindent Os testes gerados foram exportados para \texttt{Projeto1/SpotifyUM/src/test/java} e executados com a suite principal do projeto, totalizando \textbf{184 testes} na execução final validada.

\section{Projeto 2 — SpotifUM (Gradle)}
\subsection{Arquitetura e Decisões}
\subsubsection{Visão geral de classes}
A solução organiza-se em três blocos principais:
\begin{itemize}
    \item \textbf{Domínio} (módulos \texttt{model.*}): entidades de negócio como \texttt{User}, \texttt{Song}, \texttt{Playlist}, \texttt{Album} e planos de subscrição.
    \item \textbf{Utilitários} (módulo \texttt{util}): serialização/deserialização, parsing de dados e estatísticas.
    \item \textbf{Interação} (módulo \texttt{delegate}): camada de menus/fluxos de terminal.
\end{itemize}

No contexto ATS, o foco principal recaiu no núcleo de lógica de negócio e utilitários, por concentrarem comportamento funcional e serem menos voláteis do que a camada de interface textual.

\atsfigure{figures/uml_spotifum.png}{0.92}{Diagrama de classes (exportado da modelação do projeto)}

\subsubsection{Decisões técnicas}
As principais decisões de engenharia de teste foram:
\begin{itemize}
    \item privilegiar \textbf{testes orientados a risco} (estatísticas, adapters, parsing e invariantes de coleção);
    \item complementar \textbf{testes exemplo-a-exemplo} com \textbf{propriedades geradas} (jqwik);
    \item tratar mutação como mecanismo de \textbf{feedback} para identificar zonas frágeis da suite;
    \item automatizar execução para reduzir erro humano e garantir repetibilidade dos resultados.
\end{itemize}

\subsection{Metodologia de Teste}
\subsubsection{JUnit}
Foram reforçados testes unitários orientados a risco, com foco em:
\begin{itemize}
    \item regras de negócio de estatísticas exigidas no enunciado ATS;
    \item serialização/deserialização e adapters (pontos com baixa cobertura inicial);
    \item validação de fronteiras e entradas nulas/vazias.
\end{itemize}

Classes de teste adicionadas nesta fase ATS:
\begin{itemize}
    \item \texttt{LocalDateTimeAdapterTest}
    \item \texttt{SubscriptionPlanAdapterTest}
    \item \texttt{SongTypeAdapterTest}
    \item \texttt{JsonDataParserTest}
    \item \texttt{StatsPropertyTest}
    \item reforço em \texttt{StatsTest}
\end{itemize}

Foram também mantidos e estabilizados os testes existentes para evitar regressão funcional, incluindo a correção de expectativas em casos onde o comportamento do código era válido mas o teste assumia semântica incorreta.

\subsubsection{EvoSuite}
A integração foi implementada no Gradle e incluída na pipeline automática. Contudo, a execução encontra limitação de compatibilidade entre runtime disponível no container EvoSuite e bytecode Java 17 (class file 61), levando a falha de instrumentação nas classes alvo.

Resultado: requisito cumprido do ponto de vista de integração e execução técnica, mas com limitação externa de ferramenta para geração válida no ambiente atual.

\atsfigure{figures/evosuite_execution.png}{0.90}{Execução EvoSuite e limitação de compatibilidade de bytecode}

\subsubsection{JaCoCo}
A recolha de cobertura foi realizada com \texttt{jacocoTestReport}, incluindo métricas de linha, branch, método e classe. Os relatórios são produzidos em XML e HTML para análise agregada e detalhada por pacote/classe.

Fórmula usada para cobertura por tipo:
\[
\text{Cobertura} = \frac{\text{covered}}{\text{covered} + \text{missed}} \times 100
\]

\subsubsection{PIT}
A estratégia aplicada foi "testes dirigidos por sobreviventes": identificar mutantes e no-coverage, criar testes específicos para esses caminhos e reexecutar PIT para validar impacto.

A leitura do PIT foi dividida em duas dimensões:
\begin{itemize}
    \item \textbf{mutation score}: eficácia global da suite para detetar mutantes;
    \item \textbf{no coverage}: volume de código mutado não exercitado por testes.
\end{itemize}

\subsubsection{Property-Based Testing (jqwik)}
Foram implementadas propriedades para invariantes de comportamento, incluindo:
\begin{itemize}
    \item invariantes de \texttt{Playlist} (substituição de lista e cópia defensiva);
    \item consistência de contagem em \texttt{Stats.countPublicPlaylists} com entradas geradas automaticamente.
\end{itemize}

As propriedades foram executadas com múltiplas tentativas aleatórias por execução, aumentando a probabilidade de detetar defeitos em combinações não cobertas por testes example-based.

\subsection{Resultados}
\subsubsection{Resultados JUnit}
\begin{table}[H]
\centering
\begin{tabular}{lr}
\toprule
Métrica & Valor \\
\midrule
Total de testes executados & 141 \\
Taxa de sucesso & 100\% \\
\bottomrule
\end{tabular}
\caption{Resumo da execução JUnit}
\end{table}

\atsfigure{figures/junit_report_overview.png}{0.94}{Resumo da execução JUnit (Gradle test report)}

\subsubsection{Resultados de cobertura (JaCoCo)}
\begin{table}[H]
\centering
\begin{tabular}{lrrr}
\toprule
Tipo & Cobertas & Falhadas & Cobertura \\
\midrule
Line & 682 & 1022 & 40.02\% \\
Branch & 263 & 339 & 43.69\% \\
Method & 183 & 126 & 59.22\% \\
Class & 20 & 8 & 71.43\% \\
\bottomrule
\end{tabular}
\caption{Resumo de cobertura JaCoCo}
\end{table}

\atsfigure{figures/jacoco_overview.png}{0.94}{Cobertura global JaCoCo por tipo de contador}

\subsubsection{Resultados de mutação (PIT)}
\begin{table}[H]
\centering
\begin{tabular}{lr}
\toprule
Métrica & Valor \\
\midrule
Mutações geradas & 890 \\
Mutações mortas & 265 \\
Mutation score & 30\% \\
Mutações sobreviventes & 38 \\
Mutações sem cobertura & 587 \\
Test strength & 87\% \\
Testes executados pelo PIT & 395 \\
\bottomrule
\end{tabular}
\caption{Resumo de mutação PIT}
\end{table}

\atsfigure{figures/pit_overview.png}{0.94}{Resumo PIT: score, sobreviventes e no coverage}

\subsubsection{Leitura dos resultados}
Os resultados mostram evolução positiva do baseline, especialmente em classes utilitárias e estatísticas. Ainda assim, o volume de mutantes \texttt{NO\_COVERAGE} evidencia que o maior ganho futuro está em ampliar testes para caminhos ainda não exercitados, principalmente na camada \texttt{delegate}.

\subsection{Análise Crítica}
\subsubsection{Pontos fortes}
\begin{itemize}
    \item melhoria mensurável de cobertura e mutation score face à iteração inicial;
    \item incremento da robustez em serialização/adapters, reduzindo risco de regressão em persistência;
    \item combinação de testes determinísticos com PBT para maior diversidade de cenários;
    \item pipeline automatizado reproduzível, reduzindo variabilidade operacional.
\end{itemize}

\subsubsection{Lacunas e risco residual}
\begin{itemize}
    \item persistem mutantes sobreviventes em lógica condicional e retornos em alguns métodos de negócio;
    \item o bloco \texttt{NO\_COVERAGE} no PIT continua elevado em módulos de interface textual;
    \item EvoSuite permanece condicionado por incompatibilidade de bytecode/runtime no ambiente atual.
\end{itemize}

\subsubsection{Refatoração e regressão}
Foi aplicada refatoração behavior-preserving em \texttt{Stats} para reduzir duplicação e melhorar legibilidade:
\begin{itemize}
    \item extração do helper \texttt{isNullOrEmpty(Map<?,?>)} para validações repetidas;
    \item normalização de nomenclatura local (\texttt{genreCounter}).
\end{itemize}

A regressão foi validada por reexecução de \texttt{test}, \texttt{jacocoTestReport} e \texttt{pitest}, sem introdução de falhas.

\subsection{Automação e Reprodutibilidade}
\begin{itemize}
    \item Comando único: \texttt{make full}
    \item Comandos por etapa: \texttt{make test}, \texttt{make coverage}, \texttt{make mutation}, \texttt{make property}, \texttt{make evosuite}
    \item Pipeline Gradle: \texttt{./gradlew atsFullPipeline --no-daemon}
\end{itemize}

\begin{lstlisting}[caption={Exemplo de pipeline end-to-end},label={lst:pipeline}]
make full

# equivalente Gradle
./gradlew atsFullPipeline --no-daemon
\end{lstlisting}

\atsfigure{figures/pipeline_console.png}{0.94}{Execução automatizada da pipeline ATS}

\subsection{Plano de melhoria para nota máxima}
Com base nas métricas atuais e no enunciado, os próximos ganhos de qualidade são objetivos e mensuráveis:
\begin{itemize}
    \item reduzir mutantes sobreviventes em métodos de negócio com testes dirigidos;
    \item reduzir mutantes \texttt{NO\_COVERAGE} em zonas críticas selecionadas da camada \texttt{delegate};
    \item capturar e incorporar no relatório os prints finais de JUnit, JaCoCo e PIT produzidos na última execução;
    \item manter validação de regressão a cada iteração com \texttt{make full}.
\end{itemize}

\section{Conclusão}
O trabalho nos Projetos 1 e 2 evoluiu de suites funcionais para uma abordagem de teste orientada a qualidade mensurável. A combinação de reforço manual de testes (JUnit) e, no Projeto 2, de ferramentas de cobertura (JaCoCo), mutação (PIT), property-based testing (jqwik) e automação permitiu:
\begin{itemize}
    \item elevar a confiança no núcleo funcional das aplicações;
    \item identificar de forma objetiva onde os testes ainda não exercitam o sistema (especialmente via JaCoCo/PIT no Projeto 2);
    \item transformar resultados de ferramentas em ações concretas de melhoria, mantendo validação por regressão.
\end{itemize}

Em síntese, os projetos encontram-se tecnicamente sólidos para submissão, com evidência reprodutível e análise crítica alinhada ao enunciado ATS. A principal fronteira de melhoria futura está menos em adicionar testes indiscriminadamente e mais em aumentar a eficácia sobre mutantes sobreviventes e código ainda não coberto (no Projeto 2), e em expandir de forma seletiva a cobertura de módulos ainda sem testes (no Projeto 1).

\end{document}
